% Declaration sidecar for 10.1016_j.cor.2012.08.018
% ASSISTED RESOLUTION — rung (c) of staged symbol resolution (sec:resolution).
% Written by corpusbuilder.assist, model deepseek-v4-flash, 2026-08-24.
% Non-deterministically sourced; pending human confirmation.
%
% Declaration sidecar for 10.1016_j.cor.2012.08.018
% Improving the modulo simplex algorithm for large-scale periodic timetabling
%
% STUB — written by corpusbuilder.promote because no sidecar existed.
% Fill in the ? placeholders, delete the lines that do not apply, then save
% this file as corpus/declarations/10.1016_j.cor.2012.08.018.tex and re-run promote.
%
% Every symbol below was read out of the accepted formulas. Where a line
% already carries a kind, the algebra or the reviewer settled it and the
% comment above it says which; check it, do not assume it. Where three lines
% appear, decide whether the symbol is an index family, a parameter or a
% variable and delete the two that are wrong. Bibliographic lines
% (meta/name/desc/prov) are generated from the dossier — do not add them here.
%
% index NAME  ordered=0|1 cyclic=0|1
% param NAME  shape=I,J|- kind=scalar|vector|matrix|big_m|tolerance
% var   NAME  shape=I,J|- domain=continuous|non_negative|integer|binary
%              role=primary|auxiliary|slack|indicator
% obj         sense=min|max name=IDENT combination=sum|weighted_sum|lexicographic
% con   NAME  kind=linear|big_m|ordering|headway|capacity|flow_balance|modulo|...
%@ obj sense=min name=objective combination=weighted_sum :: Minimize the weighted sum of activity slacks sum_{A} ω_ij y_ij.
%
% PESP: classified as parameter during review
%@ param PESP shape=- kind=scalar :: Name of the periodic event scheduling problem (not an optimization parameter).
%
% y: classified as variable during review
%@ var y shape=A domain=continuous role=primary :: Slack variables y_ij = π_j − π_i − l_ij for activities (i,j) in A.
%
% unusable as an identifier, rename or drop: ω
% A: classified as index during review
%@ index A ordered=0 cyclic=0 :: Set of activities (edges) in the periodic event scheduling problem.
%
% Tz: classified as parameter during review
%@ param Tz shape=- kind=scalar :: Product of period T and modulo parameters z appearing in the periodicity constraints Γ(y+l) = Tz.
%
% l: classified as parameter during review
%@ param l shape=A kind=vector :: Lower bounds l_ij on activity durations for (i,j) in A.
%
% unusable as an identifier, rename or drop: Γ
% u: classified as parameter during review
%@ param u shape=A kind=vector :: Upper bounds u_ij on activity durations for (i,j) in A.
%
% R: classified as parameter during review
%@ param R shape=- kind=scalar :: Set of real numbers, indicating y_ij ∈ R.
%
% Z: classified as parameter during review
%@ param Z shape=- kind=scalar :: Set of integers, indicating z_ij ∈ Z.
%
% z: classified as variable during review
%@ var z shape=A domain=integer role=primary :: Modulo parameters z_ij for periodic constraints on non-tree arcs.
%
% Bound by a quantifier or big operator, so no declaration of their own: i, j
% (If one of these is really an index family, add a %@ index line for it.)
%
% Constraint rows (optional — they default to kind=linear):
%@ con eq_0004_a kind=linear :: Objective expression: minimize ∑_{A} ω_ij y_ij.
%@ con eq_0004_b kind=linear :: Periodicity constraints: Γ(y + l) = Tz.
%@ con eq_0004_c kind=linear :: Slack bounds: 0 ≤ y_ij ≤ u_ij − l_ij for all (i,j) ∈ A.
%@ con eq_0004_d kind=linear :: Slack variables are real-valued: y_ij ∈ R.
%@ con eq_0004_e kind=linear :: Modulo parameters on non-tree arcs are integer: z_ij ∈ Z.
